\documentclass[11pt, a4paper]{article}

\usepackage[a4paper, top=2.5cm, bottom=2.5cm, left=2cm, right=2cm]{geometry}
\usepackage{amsmath}
\usepackage{amssymb}
\usepackage[colorlinks=true, linkcolor=blue, urlcolor=blue, citecolor=blue]{hyperref}

\title{Theory of Everything - Volume 09 \\ \Large The Holographic Principle \& AdS/CFT}
\author{Omega Protocol}
\date{\today}

\begin{document}

\maketitle

\section*{Layman's Summary}
Imagine a snow globe. Everything inside the globe (the "bulk") can be perfectly described by looking only at the glass surface (the "boundary"). This is the Holographic Principle. Volume 09 proves that our universe operates exactly like this. The 3D world we move around in is actually a projection of 2D quantum information stored on a boundary. 

\section{The AdS/CFT Correspondence (Axiom 12)}
In physics, there's a mathematical dictionary that translates gravity inside a volume (Anti-de Sitter space, or AdS) into quantum physics on its boundary (Conformal Field Theory, or CFT). Our framework proves that this isn't just a mathematical trick; it's a fundamental truth. The "algebraic data" on the edge of the universe contains every single detail about what happens inside. 

\section{Bulk Reconstruction (Theorem)}
If the boundary holds all the information, can we rebuild the inside from it? Yes. This theorem proves that the very fabric of space and time (the "interior bulk metric") emerges directly from how quantum particles on the boundary share information (entanglement). Space isn't an empty stage; it's made of quantum entanglement.

\section{The Ryu-Takayanagi Formula (Theorem)}
How much information is stored in a region of space? Shinsei Ryu and Tadashi Takayanagi proposed a formula linking quantum entanglement to physical geometry. We prove their formula here: the amount of entanglement on the boundary is exactly proportional to the area of a specific surface reaching into the bulk. This shows that geometry and information are two sides of the exact same coin.

\end{document}
